\section{Architecture}\label{sec:architecture}

The architecture has one canonical store and one adapter per engine (Fig.~\ref{fig:arch}). The
store holds memory as RDF against \memcore. An adapter has two directions: \emph{write-in} turns
what the engine heard in a session into ontology items, and \emph{read-out} gives the engine
access to the store, either through MCP tools or by rendering the store into its prompt.
Engine-native memory, where an engine has it, is treated as a cache of the store, never as the
system of record.

\begin{figure}[t]
\centering
\resizebox{\linewidth}{!}{%
\begin{tikzpicture}[font=\scriptsize, node distance=4mm and 10mm,
  box/.style={draw, rounded corners=2pt, align=center, minimum height=8mm, inner sep=3pt},
  arr/.style={-{Stealth[length=1.6mm]}, semithick}]
\node[box, minimum width=15mm] (A) {Engine $A$\\(session)};
\node[box, right=of A] (W) {write-in adapter\\schema-constrained\\extraction};
\node[box, right=of W, fill=black!6] (S) {canonical store\\RDF / \memcore\\+ PROV, tombstones};
\node[box, right=of S] (R) {read-out adapter\\MCP tools (a1) or\\rendered prompt (a2)};
\node[box, right=of R, minimum width=15mm] (B) {Engine $B$\\(new session)};
\node[box, above=8mm of S, minimum width=118mm] (G) {registry: \memcore{} versions, JSON Schema,
  SHACL, CQ tests};
\draw[arr] (A) -- node[above]{turns} (W);
\draw[arr] (W) -- node[above]{items} (S);
\draw[arr] (S) -- (R);
\draw[arr] (R) -- node[above]{prompt} (B);
\draw[arr, dashed] (W.north |- G.south) -- node[right]{JSON Schema} (W.north);
\draw[arr, dashed] (G) -- node[right]{SHACL shapes, CQ tests} (S);
\end{tikzpicture}
}
\caption{Ontology-mediated memory. The registry publishes \memcore and the artifacts generated
from it; adapters use the JSON Schema to constrain write-in and the store to serve read-out.}
\label{fig:arch}
\end{figure}

\subsection{The \memcore Ontology}

\memcore (namespace \memns) is deliberately small: 14 classes (three of them enumerations of SKOS concepts) and 21 properties, and every
property carries a \texttt{skos:definition}. It reuses W3C vocabularies wherever they already
say what we need, and adds only what memory needs on top.

\paragraph{Items.} \texttt{mem:MemoryItem} is a \texttt{prov:Entity} with four content subclasses:
\texttt{Fact} (subject, attribute, value), \texttt{Preference} (subject, target, polarity, optional
context), \texttt{Relationship} (subject, related person, relation type) and \texttt{TaskState}
(subject, task, status, next step, optional due date). \texttt{ToolDefinition} records a tool an
agent was given, independent of any engine's calling convention. Every item has a natural-language
\texttt{mem:statement}, so a reader that knows nothing of the ontology can still use it, and a
subject (\texttt{mem:about}, ranging over \texttt{foaf:Agent}), which keeps the user's facts apart
from facts about the people they mention.

\paragraph{Closed-world shape.} The fields that identify an item are \texttt{owl:FunctionalProperty}
and carry \texttt{owl:cardinality 1} restrictions on their class. These axioms are not decoration:
the registry turns functional properties into single-valued fields and cardinality restrictions
into required fields of the generated JSON Schema, which then constrains extraction
(Sect.~\ref{sec:adapters}). The three concept-valued properties (polarity, relation type, task
status) range over classes defined by \texttt{owl:oneOf} over SKOS concepts, which the registry
emits as JSON Schema \texttt{enum}, Pydantic \texttt{Literal} and SHACL \texttt{sh:in}.

\paragraph{Time and supersession.} Memory changes. \texttt{mem:supersedes}, a subproperty of
\texttt{prov:wasRevisionOf}, links a new item to the one it replaces; an item is \emph{current}
when nothing supersedes it and it has not been erased. Superseded items stay in the store, so
the history (``you moved from Lviv to Bendigo'') remains answerable. \texttt{mem:validFrom} and
\texttt{mem:validUntil} record validity time.

\paragraph{Provenance.} Every item is \texttt{prov:wasDerivedFrom} the utterance it came from,
and each utterance is \texttt{mem:inSession} a \texttt{mem:Session} (a \texttt{prov:Activity}).
Engines are \texttt{prov:SoftwareAgent}s. The pipeline, not the model, writes provenance, so
``in which session did I tell you\dots'' is answered from the record rather than from the
engine's recollection.

\paragraph{Consent and erasure.} \texttt{mem:consentScope} limits which engines may receive an
item. Erasure removes every triple about the item and leaves a \texttt{mem:Tombstone} with the
item's IRI and a \texttt{mem:Erasure} activity (a \texttt{prov:Invalidation}) with its time, but
none of its content. An auditor can therefore verify that an erasure happened, and that the
content is gone, without learning what it was.

Listing~\ref{lst:item} shows an updated fact with its provenance.

\begin{lstlisting}[basicstyle=\ttfamily\scriptsize,frame=single,float=t,
  caption={A fact that supersedes an earlier one (prefixes omitted), from the gold store of evaluation persona e01.},label={lst:item}]
ex:m002 a mem:Fact ;
    mem:about ex:user ;
    mem:statement "Wrenna moved to Bendigo." ;
    mem:validFrom "2026-07-15T12:00:00Z"^^xsd:dateTime ;
    mem:portability mem:Portable ;
    mem:attribute "home city" ;
    mem:value "Bendigo" ;
    mem:supersedes ex:m001 ;
    prov:wasDerivedFrom ex:u5_15 .
ex:u5_15 a mem:Utterance ; mem:inSession ex:s5 .
\end{lstlisting}

\subsection{Ontology as a Versioned Package}\label{sec:registry}

\memcore is published in an ontology registry (\registry) that treats ontologies like software
packages. Each version is immutable and identified by SemVer and two hashes: an integrity hash of
the serialised package and a semantic hash of its axioms, so a change of formatting is
distinguishable from a change of meaning. The registry enforces the SemVer bump a change
requires, records dependencies at exact versions (here PROV-O, FOAF, DCMI Terms, OWL-Time and
SKOS), and runs the package's competency questions~\cite{gruninger1995} as tests on publication.

\memcore ships eleven competency questions as SPARQL queries with expected results on an example
store, among them: \emph{which preferences are current}, \emph{which facts were superseded and by
what}, \emph{in which session was an item learnt}, \emph{which items may be shared with engine
X}, and \emph{is every erased item absent apart from its tombstone}. Version 0.4.0, the version
used in this paper, passes all of them and covers 35 of 35 declared terms.

For the adapters, the important artifact is the generated JSON Schema: a closed-world projection
of each class with types from ranges, required fields from cardinalities, single values from
functional properties, enumerations from \texttt{owl:oneOf}, and descriptions from
\texttt{skos:definition}. The adapters never hand-write a schema; they pin a registry version and
derive everything from it.

\subsection{Adapters}\label{sec:adapters}

\paragraph{Write-in} runs after every session. The engine receives the session, a listing of the
current store (so it can link supersessions and act on erasure requests), and a field guide built
from the schema's descriptions. Its reply is constrained by a JSON Schema derived from the
registry schema: Ollama's structured output (or a hosted provider's) turns the schema into a
grammar, so the reply is syntactically valid by construction. The adapter removes from the model's view the fields the
pipeline owns (identifiers, provenance, validity, consent, portability class), maps names to
person IRIs, converts enumeration labels to concept IRIs, validates every item against the
registry schema, and writes it. Supersession and erasure references must point to existing
items; invalid references are counted, not repaired.

\paragraph{Read-out} has two forms. With \emph{tools}, the engine queries an MCP server with
four tools: \texttt{search\_memory}, \texttt{list\_people}, \texttt{memories\_about} (by name or
relation) and \texttt{memory\_history} (including superseded items). With \emph{rendering}, the
current items are written into the prompt, grouped by person and annotated with that person's
relation to the user. Both forms show only current, non-erased items, apart from the history tool.

\paragraph{Why constraint is not enough.} An early version of \memcore made every property an
optional array. Its generated schema was satisfiable by almost anything, and on a development
excerpt all three engines produced output that validated and was nonsense: facts with no
attribute, person IRIs used as item identifiers, invented classes. Declaring functional
properties, cardinalities and definitions in the ontology, so that they flow through the registry
into the schema and the prompt, is what made extraction usable (development observation; see the
repository's findings log). Schema-valid is not the same as meaningful; the ontology's
closed-world shape carries the meaning into the grammar.
